Before arriving at the tensorised expression for the angular power spectra in Section~\ref{sec:5}, we first express each factor inside the Limber integral (Equation~\ref{eq:2.1}) in piecewise-linear form on the discrete redshift grid defined in Section~\ref{sec:3.3}. This section presents the linearisation of the three-dimensional power spectrum (Section~\ref{sec:4.1}), the number-density kernel (Section~\ref{sec:4.2}), and the lensing convergence kernel (Section~\ref{sec:4.3}), together with their validation against \texttt{CCL}. The central result is that both kernel functions can be decomposed as linear combinations of the discretised tomographic redshift distribution values \(\phi^{a}_{i}\), with coefficients that depend only on cosmology and the redshift grid. Under the spatially flat cosmology adopted in Section~\ref{sec:3.1}, we have \(f_K(\chi) = \chi\) throughout; the extension to curved geometries, which amounts to linearising \(f_K(\chi)\) itself, is deferred to future work.

\subsection{Three-dimensional power spectrum} \label{sec:4.1}

Within each redshift interval \([\chi_n, \chi_{n+1}]\), the three-dimensional matter power spectrum evaluated at the Limber wavenumber \(k = (\ell + 1/2)/\chi\) is represented by a piecewise-linear interpolation between the values at the interval endpoints,
\begin{equation} \label{eq:4.1}
    P_{uv}\!\left[\frac{\ell + 1/2}{\chi},\, z(\chi)\right] = \frac{\chi_{n+1} - \chi}{\chi_{n+1} - \chi_n}\,P_{uv}^{\ell,n} + \frac{\chi - \chi_n}{\chi_{n+1} - \chi_n}\,P_{uv}^{\ell,n+1},
\end{equation}
where \(P_{uv}^{\ell,n} \equiv P_{uv}[(\ell + 1/2)/\chi_n,\, z(\chi_n)]\). This is the simplest interpolation that preserves continuity across interval boundaries and, crucially, yields integrands that can be evaluated in closed form when combined with the linearised kernel functions derived below.

Figure~\ref{fig:2} compares this piecewise-linear representation against direct evaluation by \texttt{CCL} for four representative multipoles, \(\ell = 10\), \(100\), \(1000\), and \(10\,000\). In each panel, the upper sub-panel shows the power spectrum as a function of comoving distance, while the lower sub-panel displays the ratio of each interpolation scheme to the \texttt{CCL} baseline. The \texttt{LimberCloud} piecewise-linear scheme agrees with \texttt{CCL} to better than one per cent over the comoving-distance range \(\chi \gtrsim 200\,\mathrm{Mpc}\) at all multipoles tested. Deviations exceeding a few per cent are confined to very small comoving distances (\(\chi \lesssim 200\,\mathrm{Mpc}\)), where the mapping \(k = (\ell+1/2)/\chi\) pushes the wavenumber into the deeply non-linear regime and the power spectrum varies most rapidly with distance. The first-, second-, and third-order numerical interpolation schemes reduce these residuals marginally but would break the analytic integrability on which the tensorised reformulation in Section~\ref{sec:5} relies. Since the kernel functions are negligibly small at these extreme distances for all tomographic bins considered, the impact on the final angular power spectra is negligible (see Section~\ref{sec:5}).

\begin{figure*}[!htbp]
    \centering
    \includegraphics[width=\textwidth]{Figures/POWER.pdf}
    \caption{Piecewise-linear interpolation of the non-linear matter power spectrum \(P_{\delta\delta}[(\ell+1/2)/\chi,\, z(\chi)]\) as a function of comoving distance for four representative multipoles (\(\ell = 10\), \(100\), \(1000\), and \(10\,000\); top-left to bottom-right). In each panel, the upper sub-panel compares the power spectrum evaluated by \texttt{CCL} (solid) with four interpolation schemes identified in the legend: first-order numerical interpolation, second-order numerical interpolation, third-order numerical interpolation, and the \texttt{LimberCloud} piecewise-linear scheme. The lower sub-panel shows the ratio of each scheme to \texttt{CCL}; dashed lines mark the \(\pm 1\%\) level.}
    \label{fig:2}
\end{figure*}

\subsection{Number-density kernel} \label{sec:4.2}

The number-density kernel \(W_\phi^a(\chi)\) is determined entirely by the tomographic redshift distribution. Following Equation~\ref{eq:2.3}, the distribution in comoving-distance space is related to its redshift-space counterpart by
\begin{equation} \label{eq:4.2}
    \phi^{a}_{n} \equiv \phi^a(\chi_n) = \frac{H(z_n)}{c}\,\phi^{a}(z_n),
\end{equation}
where \(n = 0, 1, \ldots, N_\mathrm{Z}\) indexes the \(N_\mathrm{Z} + 1\) grid points of the uniform redshift grid defined in Section~\ref{sec:3.3}, spanning \(N_\mathrm{Z} = 350\) intervals. Between adjacent grid points, the distribution is reconstructed by linear interpolation,
\begin{equation} \label{eq:4.3}
    \phi^{a}(\chi) = \frac{\chi_{n+1} - \chi}{\chi_{n+1} - \chi_n}\,\phi^{a}_{n} + \frac{\chi - \chi_n}{\chi_{n+1} - \chi_n}\,\phi^{a}_{n+1}, \qquad \chi_n \le \chi < \chi_{n+1},
\end{equation}
with the convention that \(\phi^{a}(\chi) = 0\) for \(\chi < \chi_0\) and \(\chi > \chi_{N_\mathrm{Z}}\).

Since the number-density kernel coincides with the redshift distribution itself, it takes the same piecewise-linear form within each interval. We denote by \(W_\phi^{a,n}(\chi)\) the restriction of the kernel to the \(n\)-th comoving-distance interval \([\chi_n, \chi_{n+1}]\); this superscript convention applies to all interval-restricted quantities throughout this work. The restricted kernel can be written as a weighted sum over all grid-point values,
\begin{equation} \label{eq:4.4}
    W_\phi^{a,n}(\chi) = \sum_{i=0}^{N_\mathrm{Z}} r_\phi^{i,n}(\chi)\,\phi^{a}_{i},
\end{equation}
where the interpolation coefficients \(r_\phi^{i,n}(\chi)\) are non-zero only for \(i = n\) and \(i = n+1\):
\begin{equation} \label{eq:4.5}
    r_\phi^{i,n}(\chi) =
    \begin{cases}
        \dfrac{\chi_{n+1} - \chi}{\chi_{n+1} - \chi_n}, & i = n, \\[8pt]
        \dfrac{\chi - \chi_n}{\chi_{n+1} - \chi_n}, & i = n+1, \\[8pt]
        0, & \text{otherwise}.
    \end{cases}
\end{equation}
The sparsity of these coefficients---only two non-zero entries per interval---reflects the local support of the piecewise-linear basis and is a key ingredient enabling the efficient tensor contraction in Section~\ref{sec:5}.

Figure~\ref{fig:3} validates the piecewise-linear reconstruction of the number-density kernel against \texttt{CCL} for all five tomographic bins of the LSST Y1 survey scenario. Each panel presents a single bin (\(a = 1\) to \(5\)), with the upper sub-panel comparing the kernel evaluated by \texttt{CCL} against four interpolation schemes---first-, second-, and third-order numerical interpolation, together with the \texttt{LimberCloud} piecewise-linear reconstruction---and the lower sub-panel displaying the ratio to \texttt{CCL}. The \texttt{LimberCloud} scheme reproduces the number-density kernels to sub-per-cent accuracy across the full comoving-distance range where the kernel amplitude is non-negligible. The higher-order numerical interpolation schemes exhibit larger spurious oscillations in the ratio panels, confirming that the piecewise-linear approach offers a favourable trade-off between accuracy and analytic tractability.

\begin{figure*}[!htbp]
    \centering
    \includegraphics[width=\textwidth]{Figures/PHI_Y1.pdf}
    \caption{Piecewise-linear reconstruction of the number-density kernel \(W_\phi^a(\chi)\) for all five tomographic bins (\(a = 1\) to \(5\)) of the LSST Y1 survey scenario, arranged in a \(2 \times 3\) grid with the shared legend in the bottom-right position. In each panel, the upper sub-panel compares the kernel evaluated by \texttt{CCL} (solid) with four interpolation schemes as identified in the legend. The lower sub-panel shows the ratio of each scheme to \texttt{CCL}; dashed lines mark the \(\pm 1\%\) level.}
    \label{fig:3}
\end{figure*}

\subsection{Lensing convergence kernel} \label{sec:4.3}

The lensing convergence kernel \(W_\kappa^a(\chi)\) defined in Equation~\ref{eq:2.4} involves a line-of-sight integral over the source distribution weighted by the lensing efficiency factor \((\chi' - \chi)/\chi'\). Under the flat-sky assumption, this reads
\begin{equation} \label{eq:4.6}
    W_\kappa^{a}(\chi) = \frac{3\,\Omega_{\mathrm{m}} H_0^2}{2\,c^2}\,\frac{\chi}{a(\chi)} \int_{\chi}^{\chi_{\mathrm{H}}} \mathrm{d}\chi'\;\phi^{a}(\chi')\,\frac{\chi' - \chi}{\chi'}.
\end{equation}
When \(\chi\) lies in the \(n\)-th interval \([\chi_n, \chi_{n+1}]\), the integral is split into sub-intervals aligned with the redshift grid,
\begin{equation} \label{eq:4.7}
    W_\kappa^{a,n}(\chi) = \frac{3\,\Omega_{\mathrm{m}} H_0^2}{2\,c^2}\,\frac{\chi}{a(\chi)} \left[\int_{\chi}^{\chi_{n+1}} + \sum_{j=n+1}^{N_\mathrm{Z}-1}\int_{\chi_j}^{\chi_{j+1}}\right] \mathrm{d}\chi'\;\phi^{a}(\chi')\,\frac{\chi' - \chi}{\chi'}.
\end{equation}
Substituting the piecewise-linear representation of \(\phi^{a}(\chi')\) from Equation~\ref{eq:4.3} into each sub-integral and collecting all contributions, the kernel can be expressed as a linear combination of the discretised distribution values,
\begin{equation} \label{eq:4.8}
    W_\kappa^{a,n}(\chi) = \frac{3\,\Omega_{\mathrm{m}} H_0^2}{2\,c^2}\,\frac{\chi}{a(\chi)}\sum_{i=0}^{N_\mathrm{Z}} s_\kappa^{i,n}(\chi)\,\phi^{a}_{i}.
\end{equation}
Unlike the number-density case, the coefficients \(s_\kappa^{i,n}(\chi)\) are non-zero for all \(i \ge n\), reflecting the non-local nature of the lensing kernel: matter at comoving distance \(\chi\) receives lensing contributions from all source galaxies at \(\chi' > \chi\). Each sub-integral involves a product of two linear functions of \(\chi'\) divided by \(\chi'\), which admits a closed-form evaluation involving logarithmic terms. The coefficients vanish for \(i < n\) and are given for each remaining case by
\begin{align}
    s_\kappa^{i=n,\,n}(\chi) &= \frac{\chi_{n+1}^2 - \chi^2}{2(\chi_{n+1} - \chi_n)} - \frac{\chi\,\chi_{n+1}}{\chi_{n+1} - \chi_n}\ln\frac{\chi_{n+1}}{\chi}, \label{eq:4.9a} \\[6pt]
    s_\kappa^{i=n+1,\,n}(\chi) &= \frac{\chi_{n+2} - \chi_{n+1}}{2} - \frac{\chi\,\chi_{n+2}}{\chi_{n+2} - \chi_{n+1}}\ln\frac{\chi_{n+2}}{\chi_{n+1}} \nonumber\\
    &\quad + \frac{\chi_{n+1}^2 - 2\chi_n\chi_{n+1} + \chi^2}{2(\chi_{n+1} - \chi_n)} + \frac{\chi\,\chi_n}{\chi_{n+1} - \chi_n}\ln\frac{\chi_{n+1}}{\chi}, \label{eq:4.9b} \\[6pt]
    s_\kappa^{n+1<i<N_\mathrm{Z},\,n}(\chi) &= \frac{\chi_{i+1} - \chi_i}{2} - \frac{\chi\,\chi_{i+1}}{\chi_{i+1} - \chi_i}\ln\frac{\chi_{i+1}}{\chi_i} \nonumber\\
    &\quad + \frac{\chi_i - \chi_{i-1}}{2} + \frac{\chi\,\chi_{i-1}}{\chi_i - \chi_{i-1}}\ln\frac{\chi_i}{\chi_{i-1}}, \label{eq:4.9c} \\[6pt]
    s_\kappa^{i=N_\mathrm{Z},\,n}(\chi) &= -\chi + \frac{\chi_{N_\mathrm{Z}} - \chi_{N_\mathrm{Z}-1}}{2} + \frac{\chi\,\chi_{N_\mathrm{Z}-1}}{\chi_{N_\mathrm{Z}} - \chi_{N_\mathrm{Z}-1}}\ln\frac{\chi_{N_\mathrm{Z}}}{\chi_{N_\mathrm{Z}-1}}. \label{eq:4.9d}
\end{align}
The \(i = n\) case involves an integral starting from \(\chi\) (the evaluation point) rather than from a grid boundary, which introduces \(\chi\)-dependent terms; all other cases integrate over complete grid intervals and depend on \(\chi\) only through the lensing efficiency factor \((\chi' - \chi)/\chi'\).

The remaining factor \(1/a(\chi) = 1 + z(\chi)\) is itself linearised within each interval as
\begin{equation} \label{eq:4.10}
    1 + z(\chi) = 1 + \frac{\chi_{n+1} - \chi}{\chi_{n+1} - \chi_n}\,z_n + \frac{\chi - \chi_n}{\chi_{n+1} - \chi_n}\,z_{n+1}, \qquad \chi_n \le \chi < \chi_{n+1}.
\end{equation}
Absorbing this factor together with the cosmological prefactor into a single set of coefficients, the full lensing convergence kernel takes the same structural form as the number-density kernel,
\begin{equation} \label{eq:4.11}
    W_\kappa^{a,n}(\chi) = \sum_{i=0}^{N_\mathrm{Z}} r_\kappa^{i,n}(\chi)\,\phi^{a}_{i},
\end{equation}
where
\begin{equation} \label{eq:4.12}
    r_\kappa^{i,n}(\chi) = \frac{3\,\Omega_{\mathrm{m}} H_0^2}{2\,c^2}\,\chi\left(1 + \frac{\chi_{n+1} - \chi}{\chi_{n+1} - \chi_n}\,z_n + \frac{\chi - \chi_n}{\chi_{n+1} - \chi_n}\,z_{n+1}\right) s_\kappa^{i,n}(\chi).
\end{equation}

Figure~\ref{fig:4} validates the piecewise-linear reconstruction of the lensing convergence kernel against \texttt{CCL} for all five tomographic bins of the LSST Y1 survey scenario. The panel layout and line conventions follow Figure~\ref{fig:3}. The \texttt{LimberCloud} scheme reproduces the lensing convergence kernels with residuals remaining below one per cent throughout the comoving-distance range that dominates the Limber integral. The oscillatory residuals visible at large comoving distances in the ratio panels---most pronounced for the higher-order numerical schemes---arise where the kernel amplitude has fallen to a negligible fraction of its peak value and do not affect the final angular power spectra. The broader support of \(W_\kappa^a\) compared with \(W_\phi^a\) reflects the non-local nature of the lensing efficiency integral, which accumulates contributions from all source galaxies beyond the evaluation point.

\begin{figure*}[!htbp]
    \centering
    \includegraphics[width=\textwidth]{Figures/KAPPA_Y1.pdf}
    \caption{Piecewise-linear reconstruction of the lensing convergence kernel \(W_\kappa^a(\chi)\) for all five tomographic bins (\(a = 1\) to \(5\)) of the LSST Y1 survey scenario, arranged in a \(2 \times 3\) grid with the shared legend in the bottom-right position. In each panel, the upper sub-panel compares the kernel evaluated by \texttt{CCL} (solid) with four interpolation schemes as identified in the legend. The lower sub-panel shows the ratio of each scheme to \texttt{CCL}; dashed lines mark the \(\pm 1\%\) level.}
    \label{fig:4}
\end{figure*}

The central result of this section is that both kernel functions admit the same structural decomposition: within each redshift interval, the kernel for tomographic bin \(a\) is a linear function of the discretised distribution values \(\phi^{a}_{i}\), with cosmology- and geometry-dependent coefficients \(r_\phi^{i,n}\) or \(r_\kappa^{i,n}\) that are independent of the survey-specific redshift distributions. This clean separation is the foundation of the tensorised reformulation presented in Section~\ref{sec:5}.
